\chapter{Estado del Arte y Fundamentos Teóricos}
\label{cap:estado_del_arte_fund_teoricos}

\section{Fundamentos de la Tomografía por Emisión de Positrones (PET)}
\label{sec:estado_del_arte_fund_teoricos:fundamentos_PET}

\subsection{Principio de Detección por Coincidencia.}
\subsection{Física de la Aniquilación y Rango del Positrón.}
\subsection{Reconstrucción de la Imagen}
\subsection{Causas de Degradación de la Imagen}

\section{Física de los Isótopos "Sucios" (\yochoseis)}
\label{sec:estado_del_arte_fund_teoricos:isotopos_sucios}

\subsection{Esquema de desintegración del Itrio-86.}
\subsection{Fenómeno de Prompt Gammas (Gamas en cascada).}
\subsection{Tipos de eventos: Verdaderos, Scatter, Randoms y Coincidencias Triples.}

\section{Simulación Monte Carlo en Medicina Nuclear}
\label{sec:estestado_del_arte_fund_teoricos:simulacion}

\subsection{Introducción al método Monte Carlo.}
\subsection{El código PENELOPE y el framework penRed.}
\subsection{Simulación de fuentes complejas (módulo PENNUC).}

\section{Imagen PET Multiplexada}
\label{sec:estado_del_arte_fund_teoricos:multiplexing}

\subsection{Concepto y estado actual de la técnica.}
\subsection{Metodologías de separación de isótopos basadas en triples.}
